\chapter{กรอบการดำเนินงาน}

\section{โครงสร้าง 3 โครงการหลักตามแบบ ง8}
การดำเนินงานทั้งหมดจัดอยู่ภายใต้โครงการหลัก 3 โครงการตามแบบ ง8 ซึ่งมีบทบาทต่างกันชัดเจน
คือ ด้านเทคโนโลยี ด้านองค์ความรู้และแหล่งเรียนรู้ และด้านกำลังคน

\begin{table}[H]
\centering
\caption{โครงการหลักตามแบบ ง8 และการเบิกจ่าย ณ 4 สิงหาคม 2569}
\small
\begin{tabular}{L{6.4cm}rrrr}
\toprule
\rowcolor{headerblue}
\textbf{โครงการหลัก (ง8)} & \textbf{โครงการ} & \textbf{งบจัดสรร} & \textbf{เบิกจ่าย} & \textbf{ร้อยละ}\\
\midrule
ง8-1 ผลักดันเทคโนโลยี นวัตกรรมสู่ชุมชน ตามเป้าหมายการพัฒนาอย่างยั่งยืน & 16 & 1,988,000 & 1,064,116 & 54\%\\
\rowcolor{rowblue}
ง8-2 ขับเคลื่อนกลไกการพัฒนาองค์ความรู้ เพื่อยกระดับคุณภาพชีวิต & 14 & 1,999,010 & 1,360,534 & 68\%\\
ง8-3 พัฒนากำลังคน สร้างอาชีพ ลดความเหลื่อมล้ำ บนพื้นที่สูง & 31 & 3,999,573 & 2,035,682 & 51\%\\
\midrule
\rowcolor{headerblue}
\textbf{รวมทั้งแผนงาน} & \textbf{61} & \textbf{7,986,583} & \textbf{4,460,332} & \textbf{56\%}\\
\bottomrule
\end{tabular}
\end{table}

\section{ยุทธศาสตร์และแผนงานย่อย}
โครงการทั้ง 61 โครงการอยู่ภายใต้ 5 ยุทธศาสตร์ และ 7 แผนงานย่อย ดังนี้

\begin{table}[H]
\centering
\caption{แผนงานย่อยภายใต้กลุ่มแผนงานใต้ร่มพระบารมี ปีงบประมาณ 2569}
\small
\begin{tabular}{L{1.8cm}L{12.0cm}}
\toprule
\rowcolor{headerblue}
\textbf{แผนงาน} & \textbf{ชื่อแผนงาน}\\
\midrule
ที่ 1.1 & วิจัยและพัฒนาเทคโนโลยี นวัตกรรม ที่สร้างรายได้ พัฒนาศักยภาพชุมชน และอนุรักษ์สิ่งแวดล้อมบนพื้นที่สูง\\
\rowcolor{rowblue}
ที่ 1.2 & วิจัยและพัฒนาเพื่อการแก้ปัญหาเชิงพื้นที่ที่สำคัญอย่างมีส่วนร่วมทุกภาคส่วน\\
ที่ 2.1 & วิจัยและพัฒนาเพื่อส่งเสริมอาชีพในภาคเกษตรและนอกภาคเกษตร\\
\rowcolor{rowblue}
ที่ 2.2 & วิจัยและพัฒนาเพื่อส่งเสริมอาชีพ ภายใต้มาตรฐานอาหารปลอดภัยและเป็นมิตรต่อสิ่งแวดล้อม\\
ที่ 3.1 & ยกระดับคุณภาพชีวิตและเสริมสร้างความเข้มแข็งของชุมชนอยู่ดีมีสุขตามศาสตร์พระราชา\\
\rowcolor{rowblue}
ที่ 4.1 & สนับสนุนการพัฒนาโครงสร้างพื้นฐานเพื่อเพิ่มประสิทธิภาพการดำเนินงาน\\
ที่ 5.1 & ขับเคลื่อนกลไกการปฏิบัติงานอย่างมีส่วนร่วมสู่การพัฒนาที่ยั่งยืน\\
\bottomrule
\end{tabular}
\end{table}

\section{ทฤษฎีและกรอบแนวคิดที่นำมาใช้}
การดำเนินงานวางอยู่บนฐานคิดที่ปรากฏเป็นชื่อแผนงานและชื่อโครงการจริงในปีงบประมาณ 2569
มิใช่กรอบแนวคิดที่เพิ่มเข้ามาภายหลังเพื่อประกอบรายงาน

\begin{table}[H]
\centering
\caption{กรอบแนวคิดและการนำมาใช้จริงในแผนงาน}
\small
\begin{tabular}{L{4.2cm}L{9.6cm}}
\toprule
\rowcolor{headerblue}
\textbf{กรอบแนวคิด} & \textbf{การปรากฏจริงในแผนงาน}\\
\midrule
ศาสตร์พระราชาและปรัชญาเศรษฐกิจพอเพียง & แผนงานที่ 3.1 ยกระดับคุณภาพชีวิตชุมชนอยู่ดีมีสุขตามศาสตร์พระราชา\\
\rowcolor{rowblue}
เกษตรทฤษฎีใหม่ และ โคก หนอง นา & โครงการพัฒนาพื้นที่เกษตรทฤษฎีใหม่ตามหลักปรัชญาเศรษฐกิจพอเพียง (555,000~บาท) และโครงการพัฒนาพื้นที่ต้นแบบ โคก หนอง นา\\
เข้าใจ เข้าถึง พัฒนา & โครงการยกระดับมาตรฐานผลิตภัณฑ์ผ้าใยกัญชงวิสาหกิจชุมชนดาวม่างสู่มาตรฐาน มผช.\\
\rowcolor{rowblue}
การวิจัยแบบมีส่วนร่วม & แผนงานที่ 1.2 พัฒนากระบวนการวิจัยแบบมีส่วนร่วมในระดับพื้นที่ (600,000~บาท)\\
Re-skill Up-skill New-skill & โครงการพัฒนาฐานข้อมูลด้านวิศวกรรมและยกระดับทักษะอาชีพในพื้นที่โครงการหลวง\\
\rowcolor{rowblue}
เครือข่ายความร่วมมือ MOU และ MOA & โครงการจัดทำแผนความร่วมมือและสร้างกลไกพัฒนาเครือข่ายเชิงพื้นที่\\
Logic Model & กรอบวิเคราะห์ผลผลิต ผลลัพธ์ และผลกระทบ ที่ใช้ในบทที่ 5 ของรายงานฉบับนี้\\
\bottomrule
\end{tabular}
\end{table}

\section{ทรัพยากรที่ใช้ดำเนินงาน}
แผนงานดำเนินการโดยหน่วยงานภายในมหาวิทยาลัย 11~หน่วยงาน ครอบคลุมทั้งคณะ สถาบัน วิทยาลัย
และวิทยาเขตเชียงราย น่าน ตาก และพิษณุโลก มีหัวหน้าโครงการ 42~ท่าน
และมีกิจกรรมย่อยตามแผนปฏิบัติการรวม 268~กิจกรรม กระจายในพื้นที่ดำเนินงาน 48~พื้นที่
เช่น ศูนย์พัฒนาโครงการหลวงขุนแปะ ห้วยเสี้ยว หนองหอย ทุ่งเริง แม่แฮ ปังค่า เลอตอ
สถานีเกษตรหลวงปางดะและอ่างขาง ศูนย์พัฒนาพันธุ์พืชจักรพันธ์เพ็ญศิริ
ศูนย์ศึกษาการพัฒนาห้วยฮ่องไคร้อันเนื่องมาจากพระราชดำริ และโรงเรียนพระปริยัติธรรมในพื้นที่
